\textbf{Ablation design.}
We conduct ten runs, each varying only the reward function composition while holding all other variables constant (Section 3.2).
Every run includes the binary correctness reward; the ``+'' prefix denotes which additional rewards are added (e.g., ``+ A'' means correctness + format compliance).

\vspace{1em}
\noindent\begin{minipage}{\linewidth}
\centering
\small
\begin{tabular*}{\linewidth}{@{\extracolsep{\fill}}lccccc}
\toprule
\textbf{Run} & \textbf{A (Fmt)} & \textbf{B (Prox)} & \textbf{C (Coh)} & \textbf{D (Ent)} & \textbf{E (Num)} \\
\midrule
Baseline & -- & -- & -- & -- & -- \\
\midrule
+ A & \checkmark & -- & -- & -- & -- \\
+ B & -- & \checkmark & -- & -- & -- \\
+ A+B & \checkmark & \checkmark & -- & -- & -- \\
\midrule
+ C & -- & -- & \checkmark & -- & -- \\
+ D & -- & -- & -- & \checkmark & -- \\
+ E & -- & -- & -- & -- & \checkmark \\
+ C+D & -- & -- & \checkmark & \checkmark & -- \\
+ C+D+E & -- & -- & \checkmark & \checkmark & \checkmark \\
\midrule
Comb.\ All & \checkmark & \checkmark & \checkmark & \checkmark & \checkmark \\
\bottomrule
\end{tabular*}
\captionof{table}{Reward composition per run. All runs include the binary correctness reward (omitted from columns for clarity). ``+~X'' = correctness + reward X; letter combinations = correctness + all listed rewards.}
\end{minipage}
\vspace{1em}


\textbf{Comparison tiers.}
We organize comparisons into three tiers of increasing depth:

\begin{enumerate}[nosep]
    \item \textbf{Tier 1 --- Headline comparisons}: Comb.\ All vs.\ Baseline, and each single-reward run vs.\ Baseline. These address whether additional rewards help and which ones help most.
    \item \textbf{Tier 2 --- Interaction analysis}: Each single-reward run vs.\ combination runs. This reveals whether rewards interact constructively (synergy), destructively (interference), or independently (additive).
    \item \textbf{Tier 3 --- Error distribution}: Mistake category distributions across runs. This addresses whether rewards shift the \emph{type} of errors, not just the overall error rate.
\end{enumerate}

\textbf{Significance criteria.}
Given single-seed runs without confidence intervals, we apply the following thresholds:

\begin{itemize}[nosep]
    \item Accuracy difference $< 1\%$: treated as noise.
    \item Accuracy difference $1$--$2\%$: marginal, reported with caveat.
    \item Accuracy difference $> 2\%$: likely meaningful.
    \item Error category shift $> 5$ percentage points: meaningful structural change.
    \item Rank-order change among the top-3 error categories: meaningful shift in failure profile.
\end{itemize}
